% K17 supporting section, 2026-11-21 UTC.
% Authoring/check model: GPT-6 Astra Pro; effort: not exposed.
% Bibliography entries are supplied separately in k_groups_bibitems.tex.
% Uses the parent's theorem/sourcefact/proof environments and amsmath/amssymb.

\section{The precise consequences for \texorpdfstring{$K$}{K}-groups}
\label{sec:k-groups}

\subsection{The known comparison with Vandiver}

Let \(\ell\) be an odd prime, put
\[
 S_\ell=\operatorname{Spec}\mathbb Z[1/\ell],\qquad
 A_\ell=\operatorname{Cl}(\mathbb Z[\mu_\ell])\otimes\mathbb Z_\ell,
\]
and let \(A_\ell^{[a]}\) denote the eigenspace for the \(a\)-th power
of the Teichm\"uller character, with \(a\) taken modulo \(\ell-1\).
For a finite abelian group \(M\), write \(M\{\ell\}\) for its
\(\ell\)-primary subgroup.

\begin{sourcefact}[Known comparison: BG-KUR-15, BG-WEI-82, BG-WEI-109]
For every \(k\ge1\),
\[
 K_{4k}(\mathbb Z)\{\ell\}
 \simeq H^2(S_\ell,\mathbb Z_\ell(2k+1)),                 \tag{K1}
\]
and
\[
 K_{4k}(\mathbb Z)\{\ell\}=0
 \quad\Longleftrightarrow\quad A_\ell^{[-2k]}=0.         \tag{K2}
\]
The two groups in (K2) are cyclic simultaneously.
Consequently Vandiver's conjecture at \(\ell\) is equivalent to either
of the following conditions:
\[
\begin{split}
 &K_{4k}(\mathbb Z)\{\ell\}=0
       \quad(1\le k\le(\ell-3)/2);\\
 &K_{4k}(\mathbb Z)\{\ell\}=0
       \quad\text{for every }k\ge1.
\end{split}
\tag{K3}
\]
\end{sourcefact}

Here (K1) uses the modern integral comparison edge map, followed by
localization; it is not an unspecified factorial-normalized Chern
character. Exact locators are Weibel \cite[VI.8.2, VI.8.3,
VI.10.9]{Weibel2013} and Kurihara
\cite[Lemma~1.2, Corollary~1.5, Proposition~3.7, Remark~3.11]{Kurihara1992}.
Historically, Kurihara's Proposition~3.7 gave one direction for
\(K\)-groups; his Remark~3.11 made the converse conditional on the
then-conjectural comparison. The modern unconditional equivalence
should therefore not be attributed to an unconditional 1992
comparison theorem.

We record the short character check. Over
\(\mathbb Z[\mu_\ell,1/\ell]\) the omitted prime is principal and the
Brauer group is zero. Kummer theory identifies \(H^2(\mu_\ell)\)
with the class group modulo \(\ell\). Tensoring with
\(\mu_\ell^{\otimes2k}\), and taking invariants under the group of
order \(\ell-1\), selects the exponent \(-2k\). Thus the dimensions
modulo \(\ell\) agree. Finiteness and Nakayama give both the zero
and cyclicity assertions. This argument does not identify the full
finite groups or their orders.

The trivial eigenspace is zero: extension after norm to
\(\mathbb Q\) acts on it by \(\ell-1\), while
\(\operatorname{Cl}(\mathbb Z)=0\). The other even characters are
exactly the exponents \(-2k\) for
\(1\le k\le(\ell-3)/2\). This proves (K3), including the empty range
at \(\ell=3\). For a fixed \(k\), (K2) concerns only one component.
The periodic assertion is about \emph{vanishing}, with period
\((\ell-1)/2\) in \(k\); it is not periodicity of the full finite
groups. In particular \(k=0\), where \(K_0(\mathbb Z)=\mathbb Z\),
is excluded. As a useful coverage consequence, vanishing of
\(K_{4k}\) for \(1\le k\le R\) would imply Vandiver for
every odd \(\ell\le2R+3\) (VK-COVER).

\paragraph{Background cases that do not use K17-S17.}
All regular odd primes have zero primary part in every positive
degree \(4k\). The 2-primary part is also zero in every such degree:
specialize Rognes--Weibel \cite[Theorem~0.7, proved as
Theorem~7.11]{RognesWeibel2000} to \(\mathbb Z[1/2]\), where
\(s=1,t=u=0\), and use localization and finiteness.
Furthermore,
\[
 K_4(\mathbb Z)=0,\qquad K_8(\mathbb Z)=0.              \tag{K4}
\]
The first is Rognes \cite[Theorem~0.1, p.~268]{Rognes2000};
the second is Dutour Sikiri\'c--Elbaz-Vincent--Kupers--Martinet
\cite[Theorem~6.4]{DEKM2026}, published in the January 2026 issue.
An older table listing \(K_8\) as unknown is obsolete.
Hart--Harvey--Ong \cite[\S1, item~(2)]{HartHarveyOng2017}
verify Vandiver for every odd prime \(\ell<2^{31}\).
These facts exclude all such prime divisors of every
\(|K_{4k}(\mathbb Z)|\). Also (K2) and (K4) give vanishing whenever
\(k\equiv0,1,2\pmod{(\ell-1)/2}\); the zero residue uses the
trivial class-group component.

\subsection{The negative Selmer group is exactly the required dual}

Write
\[
 D_\ell(e)=\mathbb Q_\ell/\mathbb Z_\ell(-e),\qquad
 \mathcal S_\ell(e)=
 \ker\!\left(H^1(S_\ell,D_\ell(e))
              \longrightarrow H^1(\mathbb Q_\ell,D_\ell(e))\right).
\]
All cohomology is continuous. The notation agrees with the
rational-field negative Selmer group in K17-S17-Q.

\begin{theorem}[VK-PT and VK-LOCAL]
For odd \(\ell\) and \(k\ge1\), there is a canonical isomorphism
\[
 \bigl(K_{4k}(\mathbb Z)\{\ell\}\bigr)^\vee
       \simeq\mathcal S_\ell(2k).                     \tag{K5}
\]
With the bottom coefficient ring \(O=\mathbb Z_\ell\), the
source's group
\(\operatorname{ShaU}(\chi_\ell^{-2k};2\ell)\)
is this same group.
\end{theorem}

\begin{proof}
Put \(r=2k+1\), \(T=\mathbb Z_\ell(r)\), and
\(D=T^\vee(1)=\mathbb Q_\ell/\mathbb Z_\ell(1-r)\).
At finite level set
\(A_n=T/\ell^nT=\mathbb Z/\ell^n(r)\), with reduction
transition maps, and let
\(D_n=A_n^\vee(1)=\mathbb Z/\ell^n(1-r)\), with the
dual inclusion maps into \(D_{n+1}\).
The finite Poitou--Tate sequence and local duality
\cite[I, Theorem~4.10; II, Theorem~2.9]{MilneADT}
give the segment
\[
\begin{split}
 H^1(\mathbb Q_\ell,A_n)&\longrightarrow H^1(S_\ell,D_n)^\vee
 \longrightarrow H^2(S_\ell,A_n)\\
 &\longrightarrow H^2(\mathbb Q_\ell,A_n)
 \longrightarrow H^0(S_\ell,D_n)^\vee\longrightarrow0.
\end{split}
\tag{K6}
\]
The modified real-place terms vanish since \(\ell\) is odd.
Local duality identifies \(H^2(\mathbb Q_\ell,A_n)\) with
\(H^0(\mathbb Q_\ell,D_n)^\vee\), and the last map is dual to
restriction of invariants. The global and \(\ell\)-local
cyclotomic images are both \(\mathbb Z_\ell^\times\). Their
invariants in \(D_n\) are therefore equal, including when they are
nonzero. Hence the last map in (K6) is an isomorphism.
Exactness and duality therefore give compatible exact sequences
\[
 0\longrightarrow H^2(S_\ell,A_n)^\vee
 \longrightarrow H^1(S_\ell,D_n)
 \longrightarrow H^1(\mathbb Q_\ell,D_n).
\]
This checks the direction of the dual injection before passing
to a limit.

The finite-coefficient arithmetic cohomology groups are finite,
so their inverse systems satisfy the Mittag--Leffler condition.
In particular Jannsen's exact sequence
\cite[(3.1)]{Jannsen1988}
\[
 0\longrightarrow\varprojlim\nolimits^1 H^1(S_\ell,A_n)
 \longrightarrow H^2_{\rm cont}(S_\ell,T)
 \longrightarrow\varprojlim H^2(S_\ell,A_n)
 \longrightarrow0
\]
has zero left term. Pontryagin duality identifies the dual of
this inverse limit of finite groups with
\(\varinjlim H^2(S_\ell,A_n)^\vee\).
Filtered direct limits are exact, and continuous cohomology
with discrete torsion coefficients commutes with this filtered
union: both globally and locally,
\(\varinjlim H^1(D_n)=H^1(D)\).
Taking the direct limit of the finite exact sequences proves
\[
 H^2_{\rm cont}(S_\ell,T)^\vee
 \simeq\ker\!\left(H^1(S_\ell,D)
                   \longrightarrow H^1(\mathbb Q_\ell,D)\right).
\]
Now (K1) gives (K5). Thus neither a local \(H^2\) correction
nor an unexamined inverse-limit term remains.

It remains to compare the explicit source local conditions.
For \(q\ne\ell\), \(D_\ell(e)\) is unramified and
\[
 H^1_{\rm unr}(\mathbb Q_q,D_\ell(e))
 =D_\ell(e)/(q^{-e}-1)D_\ell(e)=0 \qquad(e>0).         \tag{K7}
\]
Indeed \(q^{-e}-1\) is a nonzero \(\ell\)-adic integer and
multiplication by it is onto the divisible group
\(\mathbb Q_\ell/\mathbb Z_\ell\).
The geometric Frobenius convention replaces it by \(q^e-1\)
and gives the same result. Thus local triviality at 2 is
automatic for a class unramified away from \(\ell\).
Local triviality at \(\ell\) implies the inertia condition there.
Conversely, the source's inertia conditions give an
unramified-away-\(\ell\) class, and its decomposition condition
at \(\ell\) makes that class lie in the displayed kernel.
The odd-primary real conditions add nothing.

Finally the recovered definitions give
\(Q_O=O\otimes_{\mathbb Z_\ell}
(\mathbb Q_\ell/\mathbb Z_\ell)\) and
\(e(O)=\operatorname{length}_O(O/\ell O)\).
The bottom ring supplied by \texttt{exists\_coef} is isomorphic
to \(\mathbb Z_\ell\), so \(Q_O=D_\ell(0)\) and \(e(O)=1\).
The scalar twist is exactly \(\chi_\ell^{-e}\).
This proves the asserted normalization and group identification.
\end{proof}

\subsection{Removing every excluded prime}

\begin{theorem}[VK-EXCEPTIONS]
If \(e>0\) is even and \(\ell\) is odd with
\((\ell-1)\mid2e\), then
\[
 K_{2e}(\mathbb Z)\{\ell\}=0.                         \tag{K8}
\]
\end{theorem}

\begin{proof}
The character \(\omega^{-e}\) has order at most two.
The trivial component is already zero. Otherwise
\(-e\equiv(\ell-1)/2\pmod{\ell-1}\).
Evenness of \(e\) forces \(\ell\equiv1\pmod4\).
Put \(L=\mathbb Q(\mu_\ell)\),
\(F=\mathbb Q(\sqrt\ell)\), and \(H=\operatorname{Gal}(L/F)\).
For \(x\in A_\ell^{[-e]}\), the norm and ideal-extension maps
satisfy
\[
 jN(x)=\sum_{h\in H}h(x)=\frac{\ell-1}{2}x.           \tag{K9}
\]

We need only the elementary estimate \(h_F<\ell\).
Put \(M=\lfloor\sqrt\ell/2\rfloor\).
Minkowski's bound \cite[Theorem~4.3, p.~70]{MilneANT}
provides an integral ideal of norm at most \(M\) in each class.
If \(a_F(n)\) is the number of ideals of norm \(n\), then
\(a_F(n)\le d(n)\): at a split prime the coefficient of
the \(a\)-th power is \(a+1\), at an inert prime it is zero or
one, and at a ramified prime it is one. Multiplicativity gives
the claim. Hence
\[
 h_F\le\sum_{n=1}^M a_F(n)
 \le\sum_{n=1}^M d(n)
 =\#\{(a,b)\in\mathbb Z_{>0}^2:ab\le M\}
 \le M^2\le\frac{\ell}{4}<\ell.                      \tag{K10}
\]
Here \(\ell\ge5\) in the nontrivial quadratic case.
Thus \(\operatorname{Cl}(F)\{\ell\}=0\), so \(N(x)=0\).
The scalar \((\ell-1)/2\) is a unit on an \(\ell\)-primary
group; (K9) gives \(x=0\). Apply (K2) with \(k=e/2\).
No integral injectivity of an ideal-extension map is assumed.
\end{proof}

\subsection{The global order and cyclicity estimates}

\begin{theorem}[VK-POLY and VK-CYCLIC-RANGE]
The recovered theorem K17-S17-Q implies absolute
\(C,A\in\mathbb N\) such that, for every \(k\ge1\),
\[
 \log|K_{4k}(\mathbb Z)|\le C(2k+1)^A.               \tag{K11}
\]
Put \(B_k=C(2k+1)^A\). For every prime \(\ell\),
\[
 v_\ell\bigl(|K_{4k}(\mathbb Z)|\bigr)
       \le\left\lfloor\frac{B_k}{\log\ell}\right\rfloor. \tag{K12}
\]
In particular \(\log\ell>B_k\) implies vanishing of the primary
part, while \(2\log\ell>B_k\) makes it either zero or
\(\mathbb Z/\ell\).
\end{theorem}

\begin{proof}
The positive even \(K\)-group is finite. Its 2-primary part is
zero by the cited theorem; the 3-primary part is zero by the
trivial-character case of (K2). All primes excluded by
\((\ell-1)\mid4k\) vanish by (K8).
Apply K17-S17-Q to the finite set of remaining prime divisors
of its order. By (K5), the normalized source lengths are
exactly
\(v_\ell(|K_{4k}(\mathbb Z)|)\). The finite weighted sum is
therefore the logarithm of the order, proving (K11).
Singleton estimates give (K12). A finite
\(\mathbb Z_\ell\)-module of length at most one is zero or
\(\mathbb Z/\ell\), proving the last assertion.
\end{proof}

For odd \(\ell\), Kurihara's cyclicity comparison also makes
\(A_\ell^{[-2k]}\) cyclic in the last range. It does not
bound that class-group component's order by \(\ell\);
the comparison controls its number of generators, not its
whole order.

\subsection{Extracting exponent seven from the source proof}

\begin{theorem}[VK-EXP7]
Keeping the source's fixed levels and using the exact
highest-weight dimensions, the bound can be sharpened to
\[
 \exists C_7\in\mathbb N\quad
 \log|K_{4k}(\mathbb Z)|\le C_7(2k)^7
 \quad\text{for every }k\ge1.                       \tag{K13}
\]
The constant has not been numerically evaluated.
\end{theorem}

\begin{proof}
We identify the representation-theoretic input and then give
the full count-to-norm argument. All automorphy, cohomology,
and representation-classification inputs remain the cited
source theorems used in K17-S17.

In \texttt{WeylC.lean}, indices start at zero and
\[
 \nu_j=j+1+H-m_{n-1-j},\qquad N_{\rm det}=n+H.
\]
The helper \texttt{ExistsScalar.lean:existsScalar} calls
\texttt{Scalar.lean:scalar}, which constructs the polynomial
module generated by a nonzero highest-weight vector of
weight \(\nu\). This module is precisely one copy of
the irreducible \(L_\nu\). To check this assertion, decompose
the finite polynomial representation into irreducibles.
Every nonzero component of the highest-weight vector has
highest weight \(\nu\); in the isotypic component
\(L_\nu\otimes W\), that vector is a highest-weight vector
in \(L_\nu\) tensored with one vector of \(W\). Its orbit
span is one copy of \(L_\nu\).
The implementation takes the span over all matrices.
It equals the invertible-matrix orbit span: any linear
functional vanishing on the latter gives a polynomial that
vanishes on the Zariski-dense invertible locus, hence on all
matrices. Thus this exact dimension applies to the actual
constructed module, not just an alternative informal choice.

The helper \texttt{ExistsTens.lean:exists\_tens} gives the
required irreducible coefficient representation \(\Xi\)
after determinant twist \(-N_{\rm det}\), with
\[
 \dim\Xi\le(\dim L_\nu)^{[F:\mathbb Q]}
\]
for the totally complex fields used here. The infinitesimal
character is the same one used by the source's cohomology
theorem. The standard highest-weight classification and Weyl
dimension formula are cited background
\cite[\S\S27.2--27.3 and \S28.1]{EtingofGL2024}:
\[
 \dim L_\nu=\prod_{0\le i<j<n}
       \frac{\nu_i-\nu_j+j-i}{j-i}.                  \tag{K14}
\]

The rank-three trace-count source uses
\(m=(e-1,-1,-e-1)\), \(H=e+1\).
It gives
\[
 \nu=(2e+3,e+4,5),\qquad
 \lambda=\nu-(e+4)(1,1,1)=(e-1,0,1-e).
\]
Formula (K14) is \(e^3\), and an imaginary quadratic
field has two complex embeddings, so \(\dim\Xi_3\le e^6\).
The rank-two branch uses \(m=(e,-e)\), \(H=e\), giving
\(\nu=(2e+1,2)\) and \(\lambda=(e-1,-e)\).
Its one-embedding dimension is \(2e\); hence
\(\dim\Xi_2\le4e^2\).

Fix \(K=\mathbb Q(\sqrt{-3})\), \(N=2\).
The proof of \texttt{cardK} counts pairwise nonisomorphic
cuspidal representations by total cohomology at
\(\operatorname{KLev}(K,3,6)\).
The theorem \texttt{CohBoundFree.lean:cohBoundFree} therefore
gives a fixed positive integer \(B_0\) such that this total
cohomology is at most \(B_0\dim\Xi_3\).
It follows, by the unchanged
\texttt{CountSys.lean:exists\_eigData} construction, that
the rank-three trace data have at most \(B_0e^6\) elements.

For the rational rank-two branch the level is
\(N|\operatorname{disc}K|=6\).
The source's quadratic-base-change argument uses six fixed
imaginary quadratic fields with pairwise coprime discriminants.
For concreteness one may take discriminants
\(-3,-7,-11,-19,-23,-31\).
At any prime \(p\), five of them avoid \(p\).
For each field choose a fixed positive integer \(B_i\)
bounding total cohomology at
\(\operatorname{KLev}(K_i,2,18)\) by \(B_i\dim\Xi_2\).
These seven cohomology constants are independent of
\(e,p,\ell\).

The proof of
\texttt{CountQuadraticOfHyp.lean:count\_quadratic\_of\_hypotheses}
then applies unchanged. Its helper
\texttt{RootSum.lean:rootSum} introduces at most \(2^n\)
choices of roots of Frobenius eigenvalues and a Boolean
factor for two possible lifts. For \(n=2\) the multiplier
is \(8\). Thus, if \(I_p\) is the five-element subset
avoiding \(p\), the rational trace-data cardinality is at most
\[
 8\sum_{i\in I_p}4B_i e^2
 \le32\left(\sum_{i=1}^6B_i\right)e^2.
\]
This indexes whole eigensystems, retaining the multiplicities
needed by the norm product. Replacing them by distinct
single-prime trace values would not justify this argument.
Both trace-conversion helpers preserve indexing cardinality.
With
\[
 D=B_0+32\sum_{i=1}^6B_i,
\]
their disjoint union therefore has at most \(De^6\) elements.

Let \(p\ge7\), \(p\nmid6\), choose \(w\mid p\) in \(K\), and put
\(q=Nw=p^f\), \(1\le f\le2\). For the trivial character
orbit, the source constructs the integer
\[
 n_p(e)=\prod_{s}
       q^e\bigl(t_s-(q^{-e}+1+q^e)\bigr)
       \in\mathbb Z\setminus\{0\}.                  \tag{K15}
\]
The data are Galois stable, each scaled factor is integral,
and the purity bound \(|t_s|\le3\) makes every factor nonzero.
The exact source theorem is
\texttt{FrobeniusTraceNormNonzeroLogBound.lean:}
\texttt{frobenius\_trace\_norm\_nonzero\_log\_bound}.
It gives
\[
\begin{split}
 \log|n_p(e)|
 &\le De^6(2e\log q+1)\\
 &\le5De^7\log p,
\end{split}
\tag{K16}
\]
since \(q\le p^2\), \(e\ge1\), and \(\log p\ge1\).
The source's unchanged \texttt{OneChar} and orbit argument
then bounds the normalized negative-Selmer length by
\(2v_\ell(n_p(e))\), for \(\ell\ne p\) in the source's
admissible range. Its other arithmetic hypotheses are
exactly those checked by \texttt{ClassNumberThree.lean:gd}.

Take \(e=2k\) and \(p=7\). All excluded primary parts are
zero by (K8) and the known 2-primary theorem.
The one additional primary part \(\ell=7\) is zero because
7 is regular. Thus every prime satisfies
\[
 v_\ell(|K_{4k}(\mathbb Z)|)\le2v_\ell(n_7(2k)).
\]
Consequently
\[
 |K_{4k}(\mathbb Z)|\mid|n_7(2k)|^2,\qquad
 \log|K_{4k}(\mathbb Z)|
 \le10D\log7\,(2k)^7<20D(2k)^7.
\]
Taking \(C_7=20D\) proves (K13).
\end{proof}

This is a mathematical sharpening of the recovered proof,
not a declaration that the separate proposition
\texttt{HeightAndVanishing/Counts.lean:CorG} has been proved.
The exact source proof witnesses, without these improvements,
give exponent \(133\): the rank-three exponent is \(36\),
the six rank-two exponents are \(16\) each, and the norm
estimate adds one. Replacing sums of exponents by maxima
already gives \(37\). The exponent \(7\) additionally uses
the fixed levels and the exact coefficient dimensions.
No claim of a new compiled Lean theorem is made.

\paragraph{An exact small-degree computation target: VK-TRACE-GCD.}
Use (K15) also at \(p=11\). Since the 11-primary part
already vanishes,
\[
 |K_{4k}(\mathbb Z)|
 \mid \gcd\bigl(|n_7(2k)|,|n_{11}(2k)|\bigr)^2.       \tag{K17}
\]
Indeed each valuation is bounded by twice each of the two
integer valuations. Here \(q_7=7\), whereas \(q_{11}=121\).
A complete computation of these trace data and their gcd
could prove \(K_{12}=0\) at \(e=6\), or \(K_{16}=0\)
at \(e=8\), if the gcd has only already-excluded prime factors.
Neither trace computation has been completed in this run.
Formula (K17) identifies the finite calculation that would
be needed, including its normalization and multiplicities.

\subsection{The full-domain Soul\'e index}

\begin{theorem}[VK-SOULE and VK-SOULE-NORM]
Let \(\ell\) be odd and \(m\ge3\) odd, and put
\[
 N_\ell=\operatorname{Gal}
   (\overline{\mathbb Q}/\mathbb Q(\mu_{\ell^\infty})).
\]
For the standard Soul\'e character \(\kappa_m\) on this
full cyclotomic kernel,
\[
 [\,\mathbb Z_\ell(m):\kappa_m(N_\ell)\,]
   =|H^2(S_\ell,\mathbb Z_\ell(m))|
   =|K_{2m-2}(\mathbb Z)\{\ell\}|.                   \tag{K18}
\]
\end{theorem}

\begin{proof}
Fix compatible roots \(\zeta_n\) of order \(\ell^n\).
The standard character is given by Kummer action on roots of
\[
 E_{m,n}=
 \prod_{a\in(\mathbb Z/\ell^n)^\times}
       (1-\zeta_n^a)^{a^{m-1}}.                    \tag{K19}
\]
Kurihara \cite[p.~233, (4)]{Kurihara1992} defines the
cyclotomic class \(c(1)\in H^1(S_\ell,\mathbb Z_\ell(m))\)
at finite level by corestriction of
\((1-\zeta_n)\otimes\zeta_n^{\otimes(m-1)}\).
Restriction--corestriction to \(N_\ell\) gives the sum
of its conjugates. Factoring out
\(\zeta_n^{\otimes(m-1)}\) produces precisely (K19).
Thus \(\operatorname{res}_{N_\ell}c(1)=\kappa_m\);
there is no factorial in this normalization.

The group \(H^1(S_\ell,\mathbb Z_\ell(m))\) is free of
rank one and \(c(1)\ne0\), as recalled on the same page.
Let \(\alpha\) be a primitive generator. Its reduction in
\(H^1(S_\ell,\mathbb F_\ell(m))\) is nonzero by the
coefficient exact sequence. Write
\(\Gamma=\mathbb Z_\ell^\times=\Delta\times\Gamma_1\),
where \(|\Delta|=\ell-1\) and \(\Gamma_1=1+\ell\mathbb Z_\ell\).
Since \(m\) is odd, its residual character on \(\Delta\)
is nontrivial. Inflation--restriction gives
\[
 H^1(\Gamma,\mathbb F_\ell(m))
 =H^1(\Gamma_1,\mathbb F_\ell(m))^\Delta=0.           \tag{K20}
\]
Here \(\Gamma\) is abelian, so \(\Delta\) acts trivially
on \(\Gamma_1\) and by the nontrivial character on the
coefficient module.

If \(\alpha|_{N_\ell}\) had image contained in
\(\ell\mathbb Z_\ell(m)\), its reduction would restrict
to zero. Inflation--restriction for the cyclotomic quotient
would put that nonzero reduction in the zero group (K20),
a contradiction. Apply this first in the maximal extension
unramified outside \(\ell\); the full group \(N_\ell\)
surjects onto its cyclotomic kernel. The image of
\(\alpha|_{N_\ell}\) is a compact, hence closed subgroup
of \(\mathbb Z_\ell(m)\). Since it contains a unit,
it is the whole module.
Writing \(c(1)=a\alpha\), \(a\ne0\), now gives
\[
 [\mathbb Z_\ell(m):\kappa_m(N_\ell)]
 =[H^1(S_\ell,\mathbb Z_\ell(m)):\mathbb Z_\ell c(1)].
\]
Kurihara's Proposition~5.1 gives an inequality; his
Remark~5.2 explicitly states equality, citing Bloch--Kato,
\S6, (6.8)--(6.9). We use that equality as a known
theorem. It identifies the last index with \(|H^2|\).
The modern comparison (K1) completes (K18).
\end{proof}

This reconciles the index formula with Ishii's precise
full-domain statement \cite[Theorem~1.1, p.~2]{IshiiKummer2024}:
every odd-degree character is nonzero; the character is
surjective when \(m\equiv1\pmod{\ell-1}\); and all odd-degree
characters are surjective exactly when Vandiver holds at
\(\ell\). For a fixed degree the condition is the one
component in (K2).

Changing the chosen radical does not change the character
on \(N_\ell\), which fixes the roots of unity.
Replacing \(1-\zeta_n^a\) by \(\zeta_n^a-1\) changes a
Kummer unit by a sign, whose odd-primary Kummer class is
zero. Changing the compatible root generator by
\(b\in\mathbb Z_\ell^\times\) changes the scalar coordinate
by the unit \(b^{-m}\), so the index is unchanged.
Half-products or real norms may introduce powers of 2,
also units here. Ichimura--Kaneko's universal-series
coordinate \cite[\S1.1, pp.~313--314]{IchimuraKaneko1989}
contains \((1-\ell^{m-1})^{-1}\kappa_m/m!\):
the first factor is a unit, whereas \(m!\) need not be.
It must not be inserted into the definition (K19).
The original Ichimura--Sakaguchi paper was bibliographically
identified but not directly read in this run; no original
theorem number from it is asserted.

The proof uses oddness of \(\ell\), especially (K20), and
does not extend the index formula to 2. Nor does it
identify the image of \(\kappa_m\) on a deeper Galois
filtration subgroup with its image on all of \(N_\ell\).
Consequently it proves no integral free-generation theorem.

\subsection{A reproducible certificate at 37}

The criterion of Buhler--Harvey
\cite[\S4, author preprint pp.~4--5]{BuhlerHarvey2011}
with the printed exponent corrected by Hart--Harvey--Ong
\cite[\S7, author preprint p.~13]{HartHarveyOng2017}
says to test, for every irregular index \(a\) at \(\ell\),
\[
 V_{\ell,a}=
 \prod_{c=1}^{(\ell-1)/2}
   (\zeta^{c/2}-\zeta^{-c/2})^{c^{\ell-1-a}}\pmod q,\qquad
 V_{\ell,a}^{(q-1)/\ell}\not\equiv1\pmod q,           \tag{K21}
\]
where \(q\equiv1\pmod\ell\) is prime, \(\zeta\) has order
\(\ell\) modulo \(q\), and the half exponents are interpreted
modulo \(\ell\).
For \(\ell=37\), the residues of
\(B_2,B_4,\ldots,B_{34}\) modulo 37 are
\[
31,16,15,16,4,17,32,22,31,15,15,17,12,29,2,0,2.
\]
Thus the only irregular index is \(a=32\).
Taking \(q=149\), \(\zeta=107\), and its square root
\(s=16\) in the order-37 subgroup, we have
\(\zeta^{19c}=s^c\), since \(2\cdot19\equiv1\pmod{37}\).
Thus the computation with factors \(s^c-s^{-c}\) implements
the corrected half-exponent formula exactly and gives
\[
 16^{37}\equiv1,\quad16\not\equiv1,\quad
 V_{37,32}=48,\quad48^4=142\not\equiv1\pmod{149}.
\]
Trial division certifies the primes and hence the root's
exact order. This proves Vandiver at 37 by the cited
criterion, and consequently
\[
 K_{4k}(\mathbb Z)\{37\}=0\quad(k\ge1).
\]
The accompanying exact-arithmetic script compares two
Bernoulli algorithms and two modular-product algorithms,
and checks the equality of every square-root factor with
the corrected half-exponent factor.
The code and the application of the criterion were read
in full after running. This is a reproducible numerical
certificate with a cited arithmetic criterion, not a new
formal proof certificate (VK-37).

\subsection{The remaining numerical question}

The known primary-part exclusions give the sufficient gate
\[
 |K_{4k}(\mathbb Z)|<2^{31}
 \quad\Longrightarrow\quad K_{4k}(\mathbb Z)=0.       \tag{K22}
\]
A nontrivial finite group has a prime divisor of its order,
and all smaller primes have been excluded. Thus either a
small enough explicit bound or the exact trace gcd in
(K17) could close a fixed-degree case. Merely knowing that
an absolute constant in (K11) or (K13) exists does not.
Indeed the source's particular coarse choice
\(C=36c_5\,2^A\) satisfies \(C\ge36>\log(2^{31})\)
whenever \(c_5\ge1\); that bound cannot meet the gate even
with its smallest positive integer constant.

For a precise logical illustration, choose a prime
\(q>2^{31}\), put \(h=(q-1)/2\), and define abstract groups
\[
 M_k=
 \begin{cases}
 \mathbb Z/q,& k\equiv3\pmod h,\\
 0,&\text{otherwise}.
 \end{cases}
\]
They have no 2-primary or below-\(2^{31}\) primary parts,
\(M_1=M_2=0\), the prescribed fixed-prime vanishing
periodicity, and
\(\log|M_k|\le\lceil\log q\rceil(2k+1)^7\), but
\(M_3\ne0\). These are only a model of the specifically
listed size and numerical patterns. They are not actual
\(K\)-groups or cyclotomic class data and give no
counterexample to Vandiver (VK-NUM-GATE).

The current primary-source review found no proof of
\(K_{12}(\mathbb Z)=0\) or \(K_{16}(\mathbb Z)=0\).
This reports the scope of the search, not an exhaustive
negative theorem about the literature.
Kupers \cite[Remark~2.6, p.~4]{Kupers2019} gives an
integral Church--Farb--Putman route to \(K_{12}\).
The subsequent result of Br\"uck--Miller--Patzt--Sroka--Wilson
\cite[Theorem~B]{BrueckEtAl2024} proves rational
codimension-two vanishing, and does not state the integral
and higher-codimension input required by that route.
Proposition~\ref{prop:steinberg-two-local} extracts the localized
integral consequence of its proof; Proposition~\ref{prop:k12-two-ranks}
thereby leaves only ranks $9$ and $10$ in the $K_{12}$ homology gate.
No new specified degree vanishing is claimed from the
existential constants extracted here.
